\label{sec:introduction}

System administrators often use the Secure Shell Protocol (SSH) to deploy and manage servers on a private network or the public internet. It also secures other network connections, such as access to Git repositories containing the source code of public or private software projects. Thus, the security of SSH, and especially the security of SSH client authentication, is of critical importance. The IETF standardizes the basic SSH protocol~\cite{rfc4250,rfc4251,rfc4252,rfc4253,rfc4254}, but includes a flexible extension mechanism. Over the years, vendors have added new algorithms, features, and countermeasures to known attacks. So, to understand the real security of SSH, it is not sufficient to analyze standard documents; an investigation of actual implementations and configurations is necessary. While the cryptographic capabilities of clients and servers can be investigated in the lab, the actual \emph{configuration} of SSH clients must be measured in the real world. For servers, this can be done through internet scans (e.g.,~\cite{provos2001scanssh,6838249,USENIX:IzhTeiDur21}).

\paragraph{SSH Clients}

To measure the configuration of SSH clients, we need a strategy that is different from that for servers. Previous work used honeypots~\cite{6624967,melese2016honeypot,Wu2020MiningTI} and observed clients connecting to them via SSH. However, such data is heavily biased toward malicious clients, while we are interested in benign users. We note that SSH clients can authenticate using passwords or digital signatures. While passwords are only sent over an encrypted SSH connection, digital signatures rely on public keys that may be available for analysis. Thus, in our work, we propose to measure SSH clients that use signature-based authentication, which consists of two critical steps: secret key generation and signature generation.

\paragraph{Large-Scale Analysis of SSH Public Keys}

If a secret key is improperly generated, its security can be compromised, given only the public key. For instance, the security of RSA relies on the difficulty of the integer factorization problem. However, if RSA keys are generated incorrectly, different moduli can share a common factor. This opens the door to a batch GCD attack, which can easily factor such vulnerable keys~\cite{bernstein2004find,USENIX:HDWH12}. Similar risks exist with (EC)DSA, where flaws in key generation can lead to severe security vulnerabilities~\cite{cve-2008-0166}. To perform a large-scale analysis of SSH keys, we retrieve SSH public keys published by their owners on GitHub, GitLab, and Launchpad. We then evaluate the distribution of signature algorithms among these keys and test the keys for known weaknesses, such as RSA keys with small factors or a vulnerability to ROCA (Return of Coppersmith's Attack)~\cite{CCS:NSSKM17}. Compared to earlier work analyzing public keys from GitHub and GitLab~\cite{cox2015,EPRINT:Bock23,barbulescu2016weakrsakeys,cryptosense2015,amiet2018} (see \autoref{subsec:longitudinal}), we consider a larger set of keys from more diverse sources, and apply all known, efficient security tests against these data sources at two points in time.

\paragraph{Lab-Evaluation of SSH Client Signatures}

The client key is used in authentication to sign, among other data, the username and a \emph{session identifier} unique to the connection (\autoref{subsec:ssh}). This signature is sent to the server over an encrypted channel and cannot be observed publicly. Thus, we analyze client signatures in the lab. SSH supports RSA, DSA, ECDSA, and EdDSA. The DSA and ECDSA algorithms require a nonce for each signature. Breitner and Heninger~\cite{FC:BreHen19} have shown that weak randomness for nonce generation can lead to a secret key recovery attack, requiring only a small number of signatures for the attacker. To avoid the attack, they recommend the use of deterministic nonces, such as described in RFC~6979~\cite{rfc6979} and recently included in FIPS~186\nobreakdashes-5~\cite{FIPS-186-5}. Using a novel black-box evaluation method, we analyze the (undocumented) methods for (EC)DSA nonce generation in SSH clients. To differentiate between random and deterministic nonces, we have to generate two signatures over the exact same payload, which is impossible in multiple SSH connections due to the use of unique client and server randoms in the session identifier. To overcome this, we let the client connect to a custom server, which triggers a standard-conforming partial authentication success message, provoking the client to authenticate multiple times using the same key over the same connection. This allows us to study the nonce generation algorithm of arbitrary SSH clients and agents over just a network connection (see \autoref{sec:detnoncetest}).

In our evaluation, we found that \putty uses a non-standard deterministic nonce scheme, named \emph{$k_\text{proto}$}. We found that this algorithm outputs highly biased nonces (the nine most significant bits are always zero) when used with the elliptic curve NIST~\pfivetwoone~\cite{nist-sp-800-186}. By applying a solution to the associated hidden number problem, we can compute the private key with a success probability $\geq 0.5$ from just~\puttySignatureCount signatures, and always from~60 signatures.

\paragraph{Dangers of SSH Agents}

\begin{figure*}
    \centering
    \includegraphics[width=\textwidth]{img/SSHClientEcosystem_LinLibertine.pdf}
    \caption{The SSH client ecosystem. (1) The user connects an SSH client to a proxy server, intending to connect to a server only accessible through the proxy. (2-4,7) The SSH protocol runs between the client and the proxy. (5-6) The signature for public key authentication is created either by the client directly or by an SSH agent. (8-10,13) A second SSH protocol runs between the proxy and the intended server. (11-12) The signature is created in the SSH agent local to the user. (14) The user can run shell commands through the two connected channels. As we will show later in \autoref{subsec:sshclientfeatures}, agent forwarding (11-12) is supported by many SSH clients, two of which enable it even by default and provide a signing oracle to SSH proxies.}
    \Description{A figure depicting the complex interactions between different components of the SSH client ecosystem. First, the SSH client is either spawned directly by command invocation or indirectly by, for example, a Git client. The client then connects to the SSH proxy (if any) and performs all steps up to the SSH connection protocol. SSH authentication can be delegated to the SSH agent or cryptographic library as a separate component. Subsequently, the SSH proxy establishes a second connection to the target SSH server. The application data is transmitted through two encrypted BPP channels---from client to proxy and proxy to server.}
    \label{fig:client-eco}
\end{figure*}

Sometimes, SSH connections need to traverse firewalls and overcome network boundaries. In this case, public-facing SSH servers (``jump servers'') can be used as intermediaries, acting as a proxy between the client and the server. There are two options: The proxy server may establish TCP port forwarding to allow for a direct connection between the client and the server. However, port forwarding can also be used to exfiltrate data from an organization's network, and is thus often disabled. Second, the client can use an SSH session to the proxy server and establish a second SSH connection between the proxy and target server from there (\autoref{fig:client-eco}). In this case, the proxy must be able to authenticate itself to the server \emph{as the client}. With signature-based client authentication, this means that the \emph{proxy must be able to sign} the authentication request with \emph{the client's private key}. To enable the proxy to do so, the remote SSH client running on the proxy server can forward the signing request to a local signing oracle, the \emph{SSH agent}. The SSH agent protocol~\cite{ietf-sshm-ssh-agent-01} provides a high-bandwidth signing oracle for arbitrary messages. Thus, users are warned (for example, in the GitHub documentation~\cite{githubdocwarning}) to enable SSH agent forwarding only if they trust the proxy server. However, our analysis revealed two SSH implementations (AbsoluteTelnet and Tectia SSH) that enable SSH agent forwarding by default, ultimately giving all SSH servers that the user connects to the ability to sign arbitrary messages with the user's private key, if not explicitly disabled.

\paragraph{Contributions}

We make the following contributions:
\begin{itemize}
    \item We provide the first longitudinal analysis of SSH client signature keys in the wild, retrieved from Git-based repository hosting services (\autoref{subsec:longitudinal}). 
    We apply a thorough sequence of security tests on the collected keys, revealing 98 broken short keys, 139 keys generated from weak randomness, and 149 keys with common or small factors (\autoref{subsec:security2}).% (\autoref{subsec:securitykeys}).
    \item We evaluate the public key upload restrictions of these hosting services, which can protect users from uploading vulnerable keys to their accounts. We show that all services allow uploading some vulnerable keys, such as short RSA keys, ECDSA keys with points at infinity, or DSA keys with \textit{p} or \textit{q} not prime, among others (\autoref{subsec:upload-restrictions}).
    \item We analyze signature generation algorithms implemented in SSH clients and agents. In particular, we present a novel method to measure an SSH client's (EC)DSA nonce generation algorithm over the network (\autoref{subsec:sshclientfeatures}).
    \item We demonstrate a novel attack on \putty's deterministic nonces for ECDSA,
    recovering a NIST~\pfivetwoone private key from only \puttySignatureCount signatures 
    (\autoref{subsec:putty}).
\end{itemize}

\paragraph{Artifacts}

We share our toolset used in this research paper on Zenodo for permanent access. This toolset includes the SSH client key scraper, evaluation scripts for data aggregation, a script to generate test keys for testing public key upload restrictions, our optimized batch GCD implementation, the source code of the SSH signature sampling and analysis tool, a benchmarking script for the \putty \pfivetwoone private key recovery attack presented in \autoref{subsec:putty}, and additional material. Due to ethical concerns, we do not plan to make the database with all scraped public keys publicly available.
%
The artifacts are available at:
\begin{quote}
    \artifactsUrl
\end{quote}

\paragraph{Ethical Considerations}
\label{sec:ethics}

We downloaded the public key data from the platforms through their official APIs, adhering to their rate limits. We obtained several factorizations for RSA moduli as a result of our checks on weak keys. To reduce the risk of a potential leak before the platform can revoke the affected keys, we performed the public key collection and analysis on a separate machine with limited access. We did not use these keys in a context other than the described vulnerability analysis; we did not recover or use any private keys to authenticate to any platform.

We disclosed our findings to the platforms, allowing them to inform the affected users, revoke vulnerable keys, and implement upload restrictions as necessary. GitLab responded positively, mitigated the report internally without providing further details, and currently plans to show a warning to users when they upload DSA or weak RSA keys, starting with release 18.3.\footnote{\url{https://gitlab.com/gitlab-org/gitlab/-/issues/432624}, accessed: 2025-07-18} GitHub channels all vulnerability reports to a bug bounty program on HackerOne, whose policy states that users are self-responsible for the content in their repositories, which, to an extent, also applies to SSH public keys. A report that we opened, despite the bug bounties' policy, was closed as out of scope. We conclude that GitHub does not plan to take any further action or refine its upload restrictions. After initial response and disclosure, we did not receive any further response from Canonical regarding public keys on Launchpad.

Although the public keys we analyzed are publicly available, we currently do not plan to publish the complete datasets. The raw datasets contain millions of public keys alongside metadata, including usernames for all existing accounts on each platform and their respective account creation date. This, among other applications, could enable account linking and activity tracking across the three platforms based on the user's public key. While we do provide artifacts that can be used to create similar datasets, we deem the resources necessary to do so significant enough to justify our decision not to publish the raw datasets, thereby thwarting low-effort attacks. Additionally, not publishing the raw datasets maintains the platform operator's ability to restrict access to a user's SSH public keys in the future. The immediate evaluation results contain sensitive information, including partial and complete key factorizations of not yet revoked public keys.

We disclosed our attack on \putty (\puttycve) in early 2024. The \putty maintainer published a fix in \putty \puttyfixedversion. We notified AbsoluteTelnet and Tectia SSH about the default setting for agent forwarding---AbsoluteTelnet intends to change the default behavior in the future.
